\section{开场：从 RLHF 走向 RLVR}

这节课是 \textbf{post-training} 系列的第二讲。主线很明确：先补完上一讲的 RLHF 讨论，再把视角切到 \textbf{reinforcement learning from verifiable rewards}，也就是在那些“答案能被快速、准确地验证”的任务上，怎么把 RL 真正做成一个可扩展的训练工具。

\begin{importantbox}{这一讲的核心目标}
\begin{enumerate}
    \item 搞清楚 RLHF 为什么会遇到过优化、校准性变差、实现复杂等问题。
    \item 搞清楚 PPO 为什么能工作，但工程上很笨重。
    \item 搞清楚 GRPO 为什么会成为更轻量的替代方案。
    \item 看三个真实系统案例：DeepSeek R1、Kimi K1.5、Qwen 3。
\end{enumerate}
\end{importantbox}

\begin{figure}[H]
\centering
\includegraphics[width=0.95\textwidth]{slides-images/slide-000.jpg}
\caption{标题页：RL from verifiable rewards}
\end{figure}

\subsection{先回顾上一讲：DPO 为什么流行}

上一讲已经把 RLHF 的基础目标讲清楚了：对成对偏好数据做优化，让模型更倾向于输出“好样本”。在这个框架里，\textbf{DPO} 之所以爆火，是因为它把原本复杂的 RL 问题，压缩成了一个很像监督学习的目标。

可以把它理解成两步：
\begin{enumerate}
    \item 在非参数化假设下，把奖励和策略写成某种闭式关系。
    \item 用偏好对比损失去更新策略，让“好回答”的概率上升，“坏回答”的概率下降。
\end{enumerate}

\begin{figure}[H]
\centering
\includegraphics[width=0.95\textwidth]{slides-images/slide-003.jpg}
\caption{DPO 的直觉：去掉 reward model，去掉 on-policy rollouts，把优化变成对好坏样本的直接梯度更新}
\end{figure}

\begin{figure}[H]
\centering
\includegraphics[width=0.95\textwidth]{slides-images/slide-004.jpg}
\caption{DPO 形式化推导：先解出 implied reward，再把它重新参数化回策略}
\end{figure}

\begin{figure}[H]
\centering
\includegraphics[width=0.95\textwidth]{slides-images/slide-005.jpg}
\caption{DPO 的关键思想：把 reward 写回 policy ratio，然后对 pairwise preference 做监督式优化}
\end{figure}

\subsection{DPO 的变体和经验性结论}

课堂里特别强调了一个经验事实：\textbf{RL 论文里的结论非常依赖具体设置}。同样是 DPO、PO 或者它们的变体，不同的 base model、不同的数据、不同的 SFT 配方，最后可能给出完全不同的结果。

这也是为什么后来出现了大量 DPO 变体。课堂里点名了两个最近特别常见的方向：
\begin{itemize}
    \item \textbf{SimPO}：不再依赖 reference policy，更像是把优化目标重新写成一个更直接的打分问题。
    \item \textbf{Length-normalized DPO}：保留 DPO 的结构，但把长度偏差显式处理掉。
\end{itemize}

\begin{figure}[H]
\centering
\includegraphics[width=0.95\textwidth]{slides-images/slide-009.jpg}
\caption{DPO 的一些后续变体：SimPO 和 length-normalized DPO}
\end{figure}

\begin{importantbox}{RLHF 里的一个重要提醒}
不要把单篇论文里的实验结果当成普适真理。RLHF 的曲线、胜率、稳定性，都高度依赖任务、模型、数据和训练细节。
\end{importantbox}

\subsection{两个必须记住的问题}

课上后半段把两个问题提得很重。

第一个是 \textbf{overoptimization}。reward model 是从人类偏好或偏好代理上学来的，模型继续朝这个代理目标优化时，容易把 reward model 本身“学坏”，但真实的人类偏好并没有同步改善。

第二个是 \textbf{calibration}。RLHF 以后，模型已经不再是一个干净的概率模型；它更像一个 policy。于是很多原本从生成建模视角默认存在的“校准性”，在这里都不能再想当然。

\begin{figure}[H]
\centering
\includegraphics[width=0.95\textwidth]{slides-images/slide-011.jpg}
\caption{RLHF 训练里常见的风险：reward 过优化之后，proxy reward 和真实偏好开始分叉}
\end{figure}

\begin{figure}[H]
\centering
\includegraphics[width=0.95\textwidth]{slides-images/slide-013.jpg}
\caption{RLHF 之后模型不再天然校准，temperature=1 时经常会更自信但不一定更可靠}
\end{figure}

